\documentclass[12pt,a4paper]{article}
\usepackage[UTF8]{ctex}
\usepackage{geometry}
\usepackage{listings}
\usepackage{xcolor}
\usepackage{hyperref}
\usepackage{booktabs}
\usepackage{float}
\usepackage{tcolorbox}
\usepackage{enumitem}
\usepackage{graphicx}
\usepackage{longtable}

\geometry{left=2.5cm,right=2.5cm,top=2.5cm,bottom=2.5cm}

\lstset{
    language=Python,
    basicstyle=\ttfamily\small,
    keywordstyle=\color{blue}\bfseries,
    commentstyle=\color{gray},
    stringstyle=\color{orange},
    numbers=left,
    numberstyle=\tiny\color{gray},
    stepnumber=1,
    numbersep=8pt,
    backgroundcolor=\color{gray!10},
    frame=single,
    frameround=fttt,
    breaklines=true,
    showstringspaces=false,
    tabsize=4,
    captionpos=b,
    xleftmargin=2em,
    xrightmargin=1em,
    aboveskip=1em,
    belowskip=1em
}

\tcbuselibrary{skins,breakable}
\newtcolorbox{knowledgebox}[1][]{
    colback=blue!5!white, colframe=blue!50!black,
    fonttitle=\bfseries, title=#1, breakable}
\newtcolorbox{practicebox}[1][]{
    colback=green!5!white, colframe=green!50!black,
    fonttitle=\bfseries, title=#1, breakable}
\newtcolorbox{tipbox}[1][]{
    colback=yellow!10!white, colframe=orange!70!black,
    fonttitle=\bfseries, title=#1, breakable}
\newtcolorbox{theorybox}[1][]{
    colback=purple!5!white, colframe=purple!70!black,
    fonttitle=\bfseries, title=#1, breakable}
\newtcolorbox{warnbox}[1][]{
    colback=red!5!white, colframe=red!70!black,
    fonttitle=\bfseries, title=#1, breakable}

\hypersetup{colorlinks=true, linkcolor=blue, urlcolor=cyan, bookmarks=true}

\title{\textbf{大数据可视化实战手册}\\[0.5em]\large 从数据到图表}
\author{竞赛培训组}
\date{\today}

\begin{document}
\maketitle
\tableofcontents
\newpage

%========================================
\section{前言：什么是数据可视化？}
%========================================

\begin{theorybox}[用天气预报理解数据可视化]
天气预报告诉你"明天温度 25 度，降水概率 80\%"。你可能觉得抽象。

但如果展示一张温度曲线图 + 降水概率柱状图，你一眼就能看出"明天下午可能下雨，温度适中"。

\textbf{数据可视化}就是：把枯燥的数字变成直观的图表，让人一眼看出数据背后的信息。
\end{theorybox}

\begin{knowledgebox}[竞赛中的可视化考什么？]
根据竞赛评分标准，可视化属于"综合报告"部分（占 30 分）。要求你：
\begin{itemize}
    \item 用图表展示分析结果（不能只有数字）
    \item 图表清晰、美观、有标题和标签
    \item 图表能支持你的分析结论
\end{itemize}

\textbf{常用工具}：
\begin{itemize}
    \item \textbf{Matplotlib}：Python 最基础的绘图库（必学）
    \item \textbf{Seaborn}：基于 Matplotlib 的美化版（更美观）
    \item \textbf{Pyecharts}：基于 ECharts 的 Python 库（交互式图表，适合报告）
\end{itemize}
\end{knowledgebox}

\subsection{环境准备}

\begin{lstlisting}[language=bash]
# 安装可视化库（在 Python 环境中）
pip install matplotlib seaborn pyecharts pandas numpy

# 验证安装
python3 -c "import matplotlib; print(matplotlib.__version__)"
\end{lstlisting}

%========================================
\section{第一章：Matplotlib 基础}
%========================================

\begin{knowledgebox}[Matplotlib 是什么？]
Matplotlib 是 Python 最基础的绑图库。虽然默认样式不够美观，但它功能全面、控制力强，是所有 Python 可视化的基础。

\textbf{类比}：Matplotlib 就像一支铅笔——最基础，但你能画出任何图形。
\end{knowledgebox}

\subsection{第一个图表：柱状图}

\begin{lstlisting}[language=Python]
import matplotlib.pyplot as plt

# 设置中文显示
plt.rcParams['font.sans-serif'] = ['SimHei']  # 或 'Microsoft YaHei'
plt.rcParams['axes.unicode_minus'] = False      # 解决负号显示问题

# 数据
categories = ['手机数码', '电脑办公', '家用电器', '服饰鞋包', '美妆护肤']
sales = [1520, 980, 760, 1340, 890]

# 画柱状图
plt.figure(figsize=(10, 6))          # 设置画布大小（宽10，高6）
plt.bar(categories, sales, color='steelblue')  # 画柱状图
plt.title('各品类销售量')             # 标题
plt.xlabel('品类')                   # X轴标签
plt.ylabel('销售量')                 # Y轴标签
plt.xticks(rotation=45)              # X轴标签旋转45度
plt.tight_layout()                   # 自动调整布局
plt.savefig('bar_chart.png', dpi=150)  # 保存图片
plt.show()                           # 显示图片
\end{lstlisting}

\subsection{折线图}

\begin{theorybox}[折线图适合什么场景？]
折线图适合展示"趋势"——数据随时间的变化。比如：每天的销售趋势、每小时的访问量变化。
\end{theorybox}

\begin{lstlisting}[language=Python]
import matplotlib.pyplot as plt
import numpy as np

plt.rcParams['font.sans-serif'] = ['SimHei']
plt.rcParams['axes.unicode_minus'] = False

# 数据：7天的销售趋势
days = ['周一', '周二', '周三', '周四', '周五', '周六', '周日']
sales = [120, 135, 128, 142, 158, 210, 195]

plt.figure(figsize=(10, 6))
plt.plot(days, sales, marker='o', color='steelblue', linewidth=2)
plt.title('一周销售趋势')
plt.xlabel('日期')
plt.ylabel('销售量')
plt.grid(True, linestyle='--', alpha=0.5)  # 添加网格线
plt.savefig('line_chart.png', dpi=150)
plt.show()
\end{lstlisting}

\subsection{饼图}

\begin{theorybox}[饼图适合什么场景？]
饼图适合展示"占比"——各部分在整体中占多少比例。比如：各品类销售占比、设备类型分布。
\end{theorybox}

\begin{lstlisting}[language=Python]
import matplotlib.pyplot as plt

plt.rcParams['font.sans-serif'] = ['SimHei']
plt.rcParams['axes.unicode_minus'] = False

# 数据
categories = ['手机数码', '电脑办公', '家用电器', '服饰鞋包', '美妆护肤']
shares = [30.5, 19.7, 15.2, 26.9, 7.7]

plt.figure(figsize=(8, 8))
plt.pie(shares, labels=categories, autopct='%.1f%%', startangle=90)
plt.title('各品类销售占比')
plt.savefig('pie_chart.png', dpi=150)
plt.show()
\end{lstlisting}

\subsection{散点图}

\begin{theorybox}[散点图适合什么场景？]
散点图适合展示"两个变量之间的关系"。比如：浏览次数和购买次数的关系、价格和销量的关系。
\end{theorybox}

\begin{lstlisting}[language=Python]
import matplotlib.pyplot as plt
import numpy as np

plt.rcParams['font.sans-serif'] = ['SimHei']
plt.rcParams['axes.unicode_minus'] = False

# 模拟数据：浏览次数 vs 购买次数
np.random.seed(42)
views = np.random.randint(100, 1000, 50)
buys = views * 0.15 + np.random.randint(-20, 20, 50)  # 正相关

plt.figure(figsize=(10, 6))
plt.scatter(views, buys, alpha=0.6, color='steelblue')
plt.title('浏览次数 vs 购买次数')
plt.xlabel('浏览次数')
plt.ylabel('购买次数')
plt.grid(True, linestyle='--', alpha=0.3)
plt.savefig('scatter_plot.png', dpi=150)
plt.show()
\end{lstlisting}

\subsection{多图组合}

\begin{lstlisting}[language=Python]
import matplotlib.pyplot as plt

plt.rcParams['font.sans-serif'] = ['SimHei']
plt.rcParams['axes.unicode_minus'] = False

# 创建 2x2 的子图布局
fig, axes = plt.subplots(2, 2, figsize=(14, 10))

# 图1：柱状图
categories = ['手机', '电脑', '家电', '服饰', '美妆']
values = [1520, 980, 760, 1340, 890]
axes[0, 0].bar(categories, values, color='steelblue')
axes[0, 0].set_title('各品类销售量')

# 图2：折线图
days = ['周一', '周二', '周三', '周四', '周五', '周六', '周日']
sales = [120, 135, 128, 142, 158, 210, 195]
axes[0, 1].plot(days, sales, marker='o', color='green')
axes[0, 1].set_title('一周销售趋势')

# 图3：饼图
axes[1, 0].pie(values, labels=categories, autopct='%.1f%%')
axes[1, 0].set_title('品类占比')

# 图4：散点图
import numpy as np
x = np.random.randn(50)
y = x * 0.8 + np.random.randn(50) * 0.3
axes[1, 1].scatter(x, y, alpha=0.6, color='red')
axes[1, 1].set_title('散点图')

plt.tight_layout()
plt.savefig('combined_charts.png', dpi=150)
plt.show()
\end{lstlisting}

%========================================
\section{第二章：Seaborn 美化}
%========================================

\begin{knowledgebox}[Seaborn 是什么？]
Seaborn 是基于 Matplotlib 的"美化版"。它默认样式更好看，代码更简洁，特别适合统计数据的可视化。

\textbf{类比}：如果 Matplotlib 是铅笔，Seaborn 就是彩色铅笔——同样的基础，但画出来更好看。
\end{knowledgebox}

\subsection{安装与基础用法}

\begin{lstlisting}[language=Python]
import seaborn as sns
import matplotlib.pyplot as plt
import pandas as pd

plt.rcParams['font.sans-serif'] = ['SimHei']
plt.rcParams['axes.unicode_minus'] = False

# 设置 Seaborn 主题
sns.set_theme(style='whitegrid')  # 白色背景+网格
\end{lstlisting}

\subsection{箱线图}

\begin{theorybox}[箱线图适合什么场景？]
箱线图适合展示"数据的分布情况"——最大值、最小值、中位数、四分位数。特别适合比较不同组的数据分布。
\end{theorybox}

\begin{lstlisting}[language=Python]
import seaborn as sns
import pandas as pd
import matplotlib.pyplot as plt

# 创建模拟数据
data = pd.DataFrame({
    '品类': ['手机'] * 100 + ['电脑'] * 100 + ['家电'] * 100,
    '价格': list(np.random.normal(3000, 800, 100)) +
            list(np.random.normal(5000, 1200, 100)) +
            list(np.random.normal(2000, 600, 100))
})

plt.figure(figsize=(10, 6))
sns.boxplot(x='品类', y='价格', data=data)
plt.title('各品类价格分布（箱线图）')
plt.savefig('boxplot.png', dpi=150)
plt.show()
\end{lstlisting}

\subsection{热力图}

\begin{theorybox}[热力图适合什么场景？]
热力图适合展示"二维矩阵数据"。比如：不同城市不同品类的销售量，用颜色深浅表示数值大小。
\end{theorybox}

\begin{lstlisting}[language=Python]
import seaborn as sns
import pandas as pd
import numpy as np
import matplotlib.pyplot as plt

# 创建模拟数据：城市 x 品类 的销售矩阵
cities = ['北京', '上海', '广州', '深圳', '杭州']
categories = ['手机', '电脑', '家电', '服饰', '美妆']

data = pd.DataFrame(
    np.random.randint(100, 1000, (5, 5)),
    index=cities, columns=categories
)

plt.figure(figsize=(10, 8))
sns.heatmap(data, annot=True, fmt='d', cmap='YlOrRd')
# annot=True 显示数字, fmt='d' 整数格式, cmap 颜色方案
plt.title('城市 x 品类 销售热力图')
plt.savefig('heatmap.png', dpi=150)
plt.show()
\end{lstlisting}

\subsection{分布图}

\begin{lstlisting}[language=Python]
import seaborn as sns
import matplotlib.pyplot as plt
import numpy as np

# 直方图 + 密度曲线
data = np.random.normal(100, 20, 1000)

plt.figure(figsize=(10, 6))
sns.histplot(data, kde=True, bins=30, color='steelblue')
# kde=True 显示密度曲线
plt.title('成绩分布直方图')
plt.xlabel('成绩')
plt.ylabel('人数')
plt.savefig('distribution.png', dpi=150)
plt.show()
\end{lstlisting}

%========================================
\section{第三章：结合大数据平台的数据可视化}
%========================================

\begin{knowledgebox}[这一章要做什么？]
前面学了怎么画图，这一章教你把大数据平台上的数据读出来，然后画图。

流程：大数据平台 SQL 查询 → 导出 CSV → Python 读取 → 画图。
\end{knowledgebox}

\subsection{步骤1：从大数据平台导出数据}

\begin{lstlisting}[language=bash]
# 【在 node1 上执行】从 Hive 导出数据为 CSV

# 导出事件类型统计
hive -e "
SELECT event_type, COUNT(*) as cnt
FROM ecommerce.user_behavior
GROUP BY event_type
ORDER BY cnt DESC
" | sed 's/\t/,/g' > /tmp/event_stats.csv

# 导出品类统计
hive -e "
SELECT category,
       SUM(CASE WHEN event_type='view' THEN 1 ELSE 0 END) as views,
       SUM(CASE WHEN event_type='buy' THEN 1 ELSE 0 END) as buys
FROM ecommerce.user_behavior
GROUP BY category
" | sed 's/\t/,/g' > /tmp/category_stats.csv

# 导出时间趋势
hive -e "
SELECT SUBSTR(timestamp,12,2) as hour, COUNT(*) as cnt
FROM ecommerce.user_behavior
GROUP BY SUBSTR(timestamp,12,2)
ORDER BY hour
" | sed 's/\t/,/g' > /tmp/hourly_stats.csv

# 下载到本地（用于画图）
# scp root@node1:/tmp/event_stats.csv ./
\end{lstlisting}

\subsection{步骤2：用 Python 读取并画图}

\begin{lstlisting}[language=Python]
import pandas as pd
import matplotlib.pyplot as plt
import seaborn as sns

plt.rcParams['font.sans-serif'] = ['SimHei']
plt.rcParams['axes.unicode_minus'] = False
sns.set_theme(style='whitegrid')

# 读取数据
event_df = pd.read_csv('/tmp/event_stats.csv', names=['event_type', 'cnt'])
category_df = pd.read_csv('/tmp/category_stats.csv', names=['category', 'views', 'buys'])
hourly_df = pd.read_csv('/tmp/hourly_stats.csv', names=['hour', 'cnt'])

# 图1：事件类型分布柱状图
fig, axes = plt.subplots(2, 2, figsize=(14, 10))

axes[0, 0].barh(event_df['event_type'], event_df['cnt'], color='steelblue')
axes[0, 0].set_title('用户行为类型分布')
axes[0, 0].set_xlabel('次数')

# 图2：品类浏览量 vs 购买量
x = range(len(category_df))
axes[0, 1].bar(x, category_df['views'], width=0.4, label='浏览', align='center')
axes[0, 1].bar([i+0.4 for i in x], category_df['buys'], width=0.4, label='购买', align='center')
axes[0, 1].set_xticks([i+0.2 for i in x])
axes[0, 1].set_xticklabels(category_df['category'], rotation=45, ha='right')
axes[0, 1].set_title('各品类浏览量 vs 购买量')
axes[0, 1].legend()

# 图3：每小时活跃趋势
axes[1, 0].plot(hourly_df['hour'], hourly_df['cnt'], marker='o', color='green')
axes[1, 0].set_title('每小时活跃趋势')
axes[1, 0].set_xlabel('小时')
axes[1, 0].set_ylabel('行为次数')

# 图4：品类饼图
axes[1, 1].pie(category_df['views'], labels=category_df['category'], autopct='%.1f%%')
axes[1, 1].set_title('各品类浏览占比')

plt.tight_layout()
plt.savefig('bigdata_dashboard.png', dpi=150)
plt.show()
\end{lstlisting}

%========================================
\section{第四章：Pyecharts 交互式图表}
%========================================

\begin{knowledgebox}[Pyecharts 是什么？]
Pyecharts 是基于百度 ECharts 的 Python 库，能生成交互式 HTML 图表（鼠标悬停显示数据、可缩放）。

\textbf{适合场景}：竞赛报告需要展示效果好、有交互的图表。
\end{knowledgebox}

\subsection{安装}

\begin{lstlisting}[language=bash]
pip install pyecharts
\end{lstlisting}

\subsection{柱状图}

\begin{lstlisting}[language=Python]
from pyecharts.charts import Bar
from pyecharts import options as opts

# 数据
categories = ['手机数码', '电脑办公', '家用电器', '服饰鞋包', '美妆护肤']
sales = [1520, 980, 760, 1340, 890]

# 创建柱状图
bar = Bar()
bar.add_xaxis(categories)
bar.add_yaxis('销售量', sales)
bar.set_global_opts(
    title_opts=opts.TitleOpts(title='各品类销售量'),
    xaxis_opts=opts.AxisOpts(axislabel_opts=opts.LabelOpts(rotate=45))
)

# 生成 HTML 文件
bar.render('bar_chart.html')
print('图表已生成：bar_chart.html（用浏览器打开）')
\end{lstlisting}

\subsection{折线图}

\begin{lstlisting}[language=Python]
from pyecharts.charts import Line
from pyecharts import options as opts

days = ['周一', '周二', '周三', '周四', '周五', '周六', '周日']
sales = [120, 135, 128, 142, 158, 210, 195]

line = Line()
line.add_xaxis(days)
line.add_yaxis('销售量', sales, is_smooth=True)
line.set_global_opts(title_opts=opts.TitleOpts(title='一周销售趋势'))

line.render('line_chart.html')
\end{lstlisting}

\subsection{饼图}

\begin{lstlisting}[language=Python]
from pyecharts.charts import Pie
from pyecharts import options as opts

categories = ['手机数码', '电脑办公', '家用电器', '服饰鞋包', '美妆护肤']
shares = [30.5, 19.7, 15.2, 26.9, 7.7]

pie = Pie()
pie.add('', [list(z) for z in zip(categories, shares)])
pie.set_global_opts(title_opts=opts.TitleOpts(title='品类占比'))
pie.set_series_opts(label_opts=opts.LabelOpts(formatter='{b}: {c}%'))

pie.render('pie_chart.html')
\end{lstlisting}

\subsection{组合图表（一个页面多张图）}

\begin{lstlisting}[language=Python]
from pyecharts.charts import Bar, Line, Pie, Grid
from pyecharts import options as opts
from pyecharts.charts import Page

# 创建多个图表
bar = Bar()
bar.add_xaxis(['手机', '电脑', '家电', '服饰', '美妆'])
bar.add_yaxis('销量', [1520, 980, 760, 1340, 890])
bar.set_global_opts(title_opts=opts.TitleOpts(title='品类销量'))

line = Line()
line.add_xaxis(['周一', '周二', '周三', '周四', '周五'])
line.add_yaxis('趋势', [120, 135, 128, 142, 158])
line.set_global_opts(title_opts=opts.TitleOpts(title='销售趋势'))

pie = Pie()
pie.add('', [('手机', 30), ('电脑', 20), ('家电', 15), ('服饰', 27), ('美妆', 8)])
pie.set_global_opts(title_opts=opts.TitleOpts(title='品类占比'))

# 组合到一个页面
page = Page()
page.add(bar, line, pie)
page.render('dashboard.html')
print('仪表盘已生成：dashboard.html')
\end{lstlisting}

%========================================
\section{第五章：竞赛报告中的图表规范}
%========================================

\begin{knowledgebox}[报告图表的黄金法则]
\begin{enumerate}
    \item \textbf{每个图表必须有标题}：标题要说明"这张图在说什么"
    \item \textbf{X轴和Y轴必须有标签}：标明"这是什么维度"和"单位是什么"
    \item \textbf{图例必须有}：如果有多个数据系列，必须有图例说明
    \item \textbf{不要过度装饰}：颜色不超过 5 种，不需要 3D 效果
    \item \textbf{图表要有结论}：在图表下方写 1-2 句话说明你从图中看到了什么
\end{enumerate}
\end{knowledgebox}

\begin{tipbox}[常见图表选择指南]
\begin{itemize}
    \item \textbf{比较}不同类别 $\rightarrow$ 柱状图
    \item \textbf{趋势}随时间变化 $\rightarrow$ 折线图
    \item \textbf{占比}各部分比例 $\rightarrow$ 饼图
    \item \textbf{分布}数据集中在哪里 $\rightarrow$ 直方图、箱线图
    \item \textbf{关系}两个变量关联 $\rightarrow$ 散点图
    \item \textbf{矩阵}两个维度的值 $\rightarrow$ 热力图
\end{itemize}
\end{tipbox}

%========================================
\section{附录：Matplotlib 常用参数速查}
%========================================

{\footnotesize
\begin{longtable}{ll}
\toprule
\textbf{参数} & \textbf{说明} \\
\midrule
figsize=(w, h) & 画布大小 \\
title='标题' & 图表标题 \\
xlabel/ylabel & 坐标轴标签 \\
color='blue' & 颜色 \\
marker='o' & 数据点标记 \\
linewidth=2 & 线宽 \\
alpha=0.6 & 透明度 \\
grid(True) & 显示网格 \\
legend() & 显示图例 \\
tight\_layout() & 自动调整间距 \\
savefig('f.png', dpi=150) & 保存图片 \\
\bottomrule
\end{longtable}
}

\end{document}
